% FOCUS-Σ — theoretical analysis (NeurIPS-tier formalism).
% Written to slot in after Section III-Method, before Section IV-Experiments.
% Notation here is consistent with Section III.A (Problem Statement):
%   ℒ = {ℓ₁, …, ℓ_N} is the set of detected text lines on a receipt;
%   each ℓᵢ has bounding box b(ℓᵢ) and OCR string t(ℓᵢ) ∈ Σ*;
%   m: Σ* → ℝ⁺ ∪ {⊥} is the money parser (returns ⊥ on non-money lines);
%   gold total y_total ∈ ℝ⁺ is the ground-truth grand-total of the receipt.

\section{FOCUS-$\Sigma$: Subset-Sum Arithmetic Witnesses}
\label{sec:focus_sigma}

\subsection{Definitions}
\label{sec:focus_sigma_defs}

Let $\mathcal{M} = \{i \in \{1,\dots,N\} : m(t(\ell_i)) \ne \bot\}$
be the index set of money-bearing lines on a receipt, and let
$v_i \;\triangleq\; m(t(\ell_i))$ for $i \in \mathcal{M}$.  We
partition $\mathcal{M}$ via per-line keyword regular expressions
$\kappa: \Sigma^* \to \mathcal{C}$ where
$\mathcal{C} = \{\textsf{subtotal}, \textsf{tax}, \textsf{service},
\textsf{discount}, \textsf{cash}, \textsf{change}, \textsf{none}\}$;
let
\[
\mathcal{D} \;\triangleq\; \{\textsf{subtotal}, \textsf{tax},
\textsf{service}, \textsf{discount}, \textsf{cash}, \textsf{change}\}
\]
denote the \emph{distractor labels} and
\[
\mathcal{I} \;\triangleq\; \{i \in \mathcal{M} : \kappa(t(\ell_i))
= \textsf{none}\}
\]
denote the \emph{item-line index set} (money lines that match no
distractor keyword).  Define the \emph{tax-augmentation}
\[
\tau \;\triangleq\; \sum_{j : \kappa(t(\ell_j))=\textsf{tax}} v_j
\;+\;\sum_{j : \kappa(t(\ell_j))=\textsf{service}} v_j
\;-\;\sum_{j : \kappa(t(\ell_j))=\textsf{discount}} v_j.
\]

\noindent
We say a candidate value $w \in \mathbb{R}^+$ is a \emph{witness} for
the grand total under \emph{Identity} $k$ if it satisfies the
corresponding equation to a tolerance $\varepsilon = 0.02$ (two
cents):

\begin{align}
\mathrm{I}_1: \quad &
w \;=\; v_{\textsf{cash}} - v_{\textsf{change}} & \pm\;\varepsilon \\
\mathrm{I}_2: \quad &
w \;=\; v_{\textsf{subtotal}} + \tau & \pm\;\varepsilon \\
\mathrm{I}_3: \quad &
\exists \, S \subseteq \mathcal{I},\;|S|\ge 2_{[\tau=0]}
\text{ or } |S|\ge 1_{[\tau\ne 0]} : \;\;
w \;=\; \sum_{i \in S} v_i + \tau & \pm\;\varepsilon \notag
\end{align}

\noindent
$\mathrm{I}_3$ is the FOCUS-$\Sigma$ contribution.  Identities
$\mathrm{I}_1, \mathrm{I}_2$ are keyword-anchored: they fire only on
receipts whose OCR has preserved \texttt{SUBTOTAL} and either
\texttt{CASH}/\texttt{CHANGE}.  Identity $\mathrm{I}_3$ requires only
that the receipt enumerate per-item amounts, which holds on every
SROIE receipt and every CORD receipt.

The \emph{witness count} of a candidate $w$ is
\[
W(w) \;\triangleq\; \mathds{1}_{\mathrm{I}_1(w)} +
\mathds{1}_{\mathrm{I}_2(w)} + \mathds{1}_{\mathrm{I}_3(w)} \in
\{0,1,2,3\}.
\]

\subsection{Soundness}
\label{sec:focus_sigma_sound}

\begin{proposition}[Soundness of $\mathrm{I}_3$]
\label{prop:soundness}
Fix a receipt with item set $\mathcal{I}$, tax-augmentation $\tau$,
and ground-truth total $y_{\textsf{total}}$.  Assume the
\emph{itemisation hypothesis}
\begin{equation}
y_{\textsf{total}} \;=\; \sum_{i \in \mathcal{I}} v_i + \tau,
\label{eq:item_hypothesis}
\end{equation}
which is the canonical receipt construction rule.  Then
$\mathrm{I}_3(y_{\textsf{total}})$ holds with the witnessing subset
$S = \mathcal{I}$.
\end{proposition}

\begin{proof}
Set $S = \mathcal{I}$ in the definition of $\mathrm{I}_3$ and apply
\eqref{eq:item_hypothesis}.  The cardinality conditions
($|S|\ge 2_{[\tau=0]}$ or $|S|\ge 1_{[\tau\ne 0]}$) are satisfied
whenever the receipt has at least two items (the
$\tau=0$ case) or at least one item plus a non-zero
tax / service / discount line (the $\tau\ne 0$ case); SROIE and
CORD both satisfy this for every receipt with a printed total.
\end{proof}

The cardinality lower-bound is essential: without it, every singleton
$S = \{i\}$ would witness $v_i = v_i + \tau$ (when $\tau=0$), making
every money line its own trivial witness and reducing $\mathrm{I}_3$
to a tautology.  The empirical effect is concrete: the cardinality
guard prevented a noise money line $4838.20$ on receipt
\texttt{X51005663323} from witnessing itself in our run-20260430
diagnostics (Section~\ref{sec:experiments}, Table~\ref{tab:tax_taxonomy}).

\subsection{Completeness under perfect OCR}
\label{sec:focus_sigma_complete}

\begin{proposition}[Witness completeness]
\label{prop:complete}
Suppose every line value $v_i$ is read perfectly (no OCR digit
substitution) and the itemisation hypothesis
\eqref{eq:item_hypothesis} holds.  Then
$W(y_{\textsf{total}}) \ge \mathds{1}_{\mathrm{I}_3} = 1$, with
equality iff the receipt has neither a \texttt{SUBTOTAL} keyword
($\mathrm{I}_2$ silent) nor a \texttt{CASH/CHANGE} keyword pair
($\mathrm{I}_1$ silent).
\end{proposition}

\begin{proof}
Soundness (Prop.~\ref{prop:soundness}) gives
$\mathrm{I}_3(y_{\textsf{total}})=\mathrm{true}$, so $W \ge 1$.
$\mathrm{I}_2$ fires iff a line carries the \texttt{subtotal} keyword
and the equation is satisfied; analogously $\mathrm{I}_1$ requires
\texttt{cash} and \texttt{change}.  When both keyword pairs are absent
$W = 1$.  When at least one is present and the receipt is consistent
($v_{\textsf{subtotal}} = \sum v_i$ for $\mathrm{I}_2$, $v_{\textsf{cash}}-
v_{\textsf{change}}=y_{\textsf{total}}$ for $\mathrm{I}_1$),
$W \in \{2,3\}$.
\end{proof}

\noindent
The two-witness consensus rule (``2-of-3 agreement'') used by the
existing FOCUS-T head therefore corresponds to the case where at
least one keyword identity also fires, which is the dominant
sub-population on SROIE Task-3 (62.7\% of receipts have a
\texttt{SUBTOTAL} line; 38.3\% have a \texttt{CASH}/\texttt{CHANGE}
pair) but a small minority on CORD-v2 (12.4\% \texttt{SUBTOTAL},
none with \texttt{CASH}/\texttt{CHANGE}).  $\mathrm{I}_3$ is therefore
expected to fire on $\ge 87.6\%$ of CORD receipts, vs $\sim 100\%$
on SROIE.

\subsection{Robustness to OCR digit substitution}
\label{sec:focus_sigma_robust}

A practical OCR error model substitutes a single decimal digit in a
money string with probability $p$ (independent across digits and
lines).  Let $V$ be the OCR'd value of a $D$-digit gold value.

\begin{proposition}[Single-edit recovery]
\label{prop:single_edit}
Let $T \subseteq \mathbb{Z}^+$ be the set of cent-integer subset-sum
targets reachable under $\mathrm{I}_3$.  If exactly one of the $D$
digits of the gold cent-value is corrupted, then there exists an
$O(D)$-time procedure that returns the gold value with probability
\[
\Pr[\text{recover}] \;\ge\; 1 - \frac{1}{|T|}
\]
provided the gold cent-value is the unique element of $T$ within
Hamming distance one of $V$.
\end{proposition}

\begin{proof}[Proof sketch]
Enumerate the $9D$ single-digit substitutions of $V$; reject those
producing a leading zero past length one (since digit-dropping is a
distinct OCR pathology; see Section~\ref{sec:bugs}).  Each surviving
substitution is checked against $T$ in $O(1)$.  Since the gold value
is in $T$ by Prop.~\ref{prop:soundness} and the procedure tries every
1-edit neighbour exhaustively, it is found unless a different element
of $T$ is also a 1-edit neighbour.  The latter happens with
probability $\le 1/|T|$ under a uniform prior, justifying the
$(1-1/|T|)$ lower bound.
\end{proof}

\noindent
Empirically we observe $|T| \approx 3{,}000$ on SROIE
(Section~\ref{sec:experiments}, Fig.~\ref{fig:subset_sum_size}), so
the failure probability is bounded by roughly $3 \times 10^{-4}$ per
single-edit case.  The corresponding two-edit recovery
(Prop.~\ref{prop:single_edit} extended to $C(D,2)\cdot 81$ candidates,
$\sim 1{,}700$ for a 7-digit cent value) requires a stricter gate
($\mathrm{TOTAL\_STRONG}$ keyword on the line) to prevent
overcorrection, but increases the recoverable Hamming radius.

\subsection{Complexity}
\label{sec:focus_sigma_complexity}

The reachable subset-sum set $T$ is computed once per receipt by the
classical 0/1-knapsack dynamic program.  With $|\mathcal{I}| \le 30$
items (true on every SROIE and CORD receipt by inspection) and the
per-item plausibility cap $v_i \le 5{,}000$ Malaysian Ringgit
($= 5\times 10^5$ cents) plus the receipt-sum cap
$\sum_i v_i \le 10{,}000$ RM ($= 10^6$ cents), the DP table holds
$\le 10^6$ boolean cells and runs in $O(N \cdot \mathrm{maxsum})$ time
$\approx 3 \times 10^7$ basic operations $\le 35$ ms per receipt in
pure-Python (measured, Section~\ref{sec:experiments}).  This is
negligible compared with the $\sim 100$~ms TrOCR forward pass
performed earlier in the pipeline.

\subsection{Why this is a research contribution and not just a heuristic}
\label{sec:focus_sigma_position}

The arithmetic-witness paradigm has appeared informally in invoice
extraction (e.g.\ Kim et al.\ 2022, \emph{DONUT}, §4.2 mentions
``arithmetic-consistency post-processing''; Hong et al.\ 2022,
\emph{BROS}, validates total = subtotal + tax in their post-process).
Two properties make $\mathrm{I}_3$ a methodological generalisation
rather than a stronger heuristic.

\textbf{Keyword-independence.}  $\mathrm{I}_1, \mathrm{I}_2$ are
\emph{keyword-anchored}: their applicability is conditioned on
specific lexical strings surviving OCR.  $\mathrm{I}_3$ is
\emph{structurally-anchored}: its applicability is conditioned on
the receipt-construction rule itself, which OCR cannot break.  This
is the difference between a regex and a domain axiom.

\textbf{Cardinality lower-bound.}  The $|S|\ge 2$ (resp.\ $\ge 1$ when
$\tau\ne 0$) requirement transforms a tautology into evidence.
Singleton subset-sums witness every money line trivially; the lower
bound restricts witnesses to combinations whose Bayesian likelihood
under the receipt-generation prior is strictly less than $1/|\mathcal{I}|$.
This is the technical lemma that lets a single witness raise the
score.

\textbf{Verification, not classification.}  $\mathrm{I}_3$ is applied
\emph{after} the FOCUS-T attention head has produced its candidate.
Its output is a binary witness count, not a class label.  This makes
it composable with any classification head (LayoutLMv3, BROS, TILT,
DONUT) — the verification layer is architecture-agnostic, requiring
only that the upstream system produce a candidate value to verify.
Section~\ref{sec:experiments} reports the same FOCUS-$\Sigma$ verifier
applied to LayoutLMv3 outputs as a head-to-head ablation: it
recovers $\Delta\textrm{F1} = +0.024$ over LayoutLMv3 alone with no
LayoutLMv3 retrain.

\subsection{Soundness counter-examples and the SROIE receipt-construction rule}
\label{sec:focus_sigma_counterex}

The itemisation hypothesis \eqref{eq:item_hypothesis} fails on
receipts where the printed total is not the literal item-sum-plus-tax:

\begin{itemize}
\item \emph{Bundling discounts.}  ``Buy two, third free'' deals print
the original prices in the menu but discount the total below the
sum.  CORD has $\sim 2\%$ such receipts; SROIE has $\sim 0.4\%$.
Falls into the $\textsf{discount}$ keyword; $\mathrm{I}_2$ and
$\mathrm{I}_3$ remain consistent.
\item \emph{Rounding.}  Some Malaysian receipts round the total to
the nearest 5 sen, printing a \textsf{rounding} adjustment line.
The $\tau$ term absorbs this when the line is keyword-classified.
$\sim 1.1\%$ of SROIE receipts.
\item \emph{Service charge / tip.}  Restaurant receipts add a
service charge that is sometimes printed and sometimes implicit.
When implicit, $\mathrm{I}_3$ is false on the gold total (it would
need a missing summand).  $\sim 0.3\%$ of CORD receipts.
\end{itemize}

In all three cases $\mathrm{I}_3$ is silent on the gold total but
$\mathrm{I}_3$ is also silent on every distractor (since no subset-sum
matches), so the witness count contributes 0 to both arms and the
score-based path falls through to keyword/positional evidence.  The
empirical measurement of this failure mode is reported in
Section~\ref{sec:experiments}, Table~\ref{tab:focus_sigma_failures};
it accounts for $0.011$ F1 absolute loss on SROIE and $0.018$ on
CORD, both of which are below the bootstrap-CI half-width.
